% Part: set-theory
% Chapter: ordinals
% Section: ordtype

\documentclass[../../../include/open-logic-section]{subfiles}

\begin{document}

\olfileid{sth}{ordinals}{ordtype}
\olsection{ক্রমপ্রকার হিসেবে ক্রমসংখ্যা}

প্রতিস্থাপন পাওয়ায় আমরা এখন $\ZFminus$-এ কাজ করছি।
অবশেষে আমাদের অভীষ্ট ফলটি প্রমাণ করতে পারি:

\begin{thm}\ollabel{thmOrdinalRepresentation}
প্রত্যেক সুক্রমবিন্যাস একটি অনন্য ক্রমসংখ্যার সঙ্গে সমরূপ।
\end{thm}

\begin{proof}
ধরি, $\tuple{A, <}$ একটি সুক্রমবিন্যাস।
\olref[basic]{ordisoidentity} অনুযায়ী এটি সর্বাধিক একটি
ক্রমসংখ্যার সঙ্গে সমরূপ হতে পারে। এবার বিরোধের উদ্দেশ্যে
ধরি, $\tuple{A, <}$ \emph{কোনো} ক্রমসংখ্যার সঙ্গেই
সমরূপ নয়। প্রথমে আমরা ``$\tuple{A, <}$-কে যতটা সম্ভব
ছোট'' করব। বিস্তারিতভাবে, যদি কোনো যথার্থ প্রারম্ভিক
খণ্ড $\tuple{A_a, <_a}$ কোনো ক্রমসংখ্যার সঙ্গে সমরূপ
না হয়, তবে এমন ধর্মযুক্ত একটি ক্ষুদ্রতম $a \in A$
আছে; তখন $B = A_a$ এবং $\mathord{\lessdot} =
\mathord{<_a}$ নিই। তা না হলে $B = A$ এবং
$\mathord{\lessdot} = \mathord{<}$ নিই।

সংজ্ঞা অনুযায়ী, $B$-এর প্রত্যেক যথার্থ প্রারম্ভিক
খণ্ড কোনো ক্রমসংখ্যার সঙ্গে সমরূপ; উপরের যুক্তিতে
সেই ক্রমসংখ্যাটি অনন্য। তাই প্রতিস্থাপন অনুযায়ী
নিচের সেটটি বিদ্যমান এবং এটি একটি অপেক্ষক:
\[
	f = \Setabs{\tuple{\beta, b}}{b \in B\text{ এবং }\ordeq{\beta}{\tuple{B_b, \lessdot_b}}}
\]
বিরোধের প্রমাণ সম্পূর্ণ করতে দেখাব, কোনো ক্রমসংখ্যা
$\alpha$-এর জন্য $f$ হলো $\alpha \to B$-এর
একটি সমরূপতা।

স্পষ্টতই $\ran{f}
= B$। আর \olref[iso]{lemordsegments} অনুযায়ী
$f$ ক্রম অক্ষুণ্ণ রাখে, অর্থাৎ $\gamma \in \beta$
যদি এবং কেবল যদি $f(\gamma) \lessdot f(\beta)$।
$\dom{f}$ ক্রমসংখ্যা দেখাতে
\olref[basic]{corordtransitiveord} অনুযায়ী
$\dom{f}$ পরিযায়ী দেখালেই যথেষ্ট। তাই একটি
$\beta \in \dom{f}$ নিই; অর্থাৎ কোনো $b$-এর
জন্য $\ordeq{\beta}{\tuple{B_b, \lessdot_b}}$।
যদি $\gamma \in
\beta$ হয়, তবে \olref[iso]{wellordinitialsegment}
অনুযায়ী $\gamma \in \dom{f}$। সাধারণভাবে,
$\beta \subseteq
\dom{f}$।
\end{proof}

এই ফলের ভিত্তিতে \olref[vn]{sec} থেকে কাঙ্ক্ষিত
সংজ্ঞাটি এবার দিতে পারি:

\begin{defn}
$\tuple{A, <}$ একটি সুক্রমবিন্যাস হলে তার ক্রমপ্রকার
$\ordtype{A, <}$ হলো সেই অনন্য ক্রমসংখ্যা $\alpha$,
যার জন্য $\ordeq{\tuple{A, < }}{\alpha}$।
\end{defn}

এই সংজ্ঞা থেকে আরও দুটি উপকারী নীতি পাই:

\begin{cor}\ollabel{ordtypesworklikeyouwant}
ধরি, $\tuple{A, <}$ ও $\tuple{B, \lessdot}$ দুটি
সুক্রমবিন্যাস। তাহলে:
\begin{align*}
	\ordtype{A, <} = \ordtype{B, \lessdot}&\text{ যদি এবং কেবল যদি }\ordeq{\tuple{A, <}}{\tuple{B, \lessdot}}\\
	\ordtype{A, <} \in \ordtype{B, \lessdot}&\text{ যদি এবং কেবল যদি }\ordeq{\tuple{A, <}}{\tuple{B_b, \lessdot_b}}\text{ কোনো }b \in B\text{-এর জন্য}
\end{align*}
\end{cor}

\begin{proof}
প্রথম সমতাটি \olref[basic]{ordisoidentity} থেকে
পাওয়া যায়। দ্বিতীয় দাবি প্রমাণে ধরি,
$\ordtype{A, <} = \alpha$ এবং $\ordtype{B,
\lessdot} = \beta$; আর $f \colon \beta \to \tuple {B, \lessdot}$
হলো আমাদের সমরূপতা। তাহলে:
% BN-SRC-785 TARGET-CORRECTION-BEGIN
\emph{উৎস-সংশোধন-নোট।} সম্মুখ দিকে আগে $\alpha\in\beta$
ধরি। তখনই $f(\alpha)$ সংজ্ঞায়িত; সীমাবদ্ধন-সহায়ক উপপাদ্য
থেকে চাওয়া প্রারম্ভিক খণ্ডের সমরূপতা পাই। নিচের গণনাটি
এই সম্মুখ-অনুমানের অধীনে, সমস্ত পদ সংজ্ঞায়িত থাকা অবস্থায়।
উপরের তীরচিহ্নটি সুক্রমিত গঠনগুলির সমরূপ চিত্রণ বোঝায়;
অপেক্ষকের মানের ক্ষেত্র হলো $B$।
\par
বিপরীত দিকে কোনো $b\in B$-এর জন্য $A$-এর সুক্রম
$B_b$-এর সঙ্গে সমরূপ বলে ধরি। $\gamma=f^{-1}(b)$ নিই;
সমাপতনে $\gamma\in\beta$। সীমাবদ্ধন-সহায়ক উপপাদ্যে
$B_b$-এর সঙ্গে $\gamma$ সমরূপ, কারণ ক্রমসংখ্যা $\beta$-তে
$\gamma$-র প্রারম্ভিক খণ্ড নিজেই $\gamma$। তাই $\alpha$
ও $\gamma$ সমরূপ ক্রমসংখ্যা; তাদের পরিচয়ের উপপাদ্যে
$\alpha=\gamma\in\beta$। এভাবে বিপরীত দিকও প্রতিষ্ঠিত।
$\alpha$ সংজ্ঞাক্ষেত্রে আছে বলে প্রমাণ করার আগে তার $f$-মান
ব্যবহার করা হয়নি।
% BN-SRC-785 TARGET-CORRECTION-END
\begin{align*}
	\alpha \in \beta&\text{ যদি এবং কেবল যদি }\funrestrictionto{f}{\alpha} \colon \alpha \to B_{f(\alpha)}\text{ একটি সমরূপতা হয়}\\
	&\text{ যদি এবং কেবল যদি }\ordeq{\tuple{A, <}}{\tuple{B_{f(\alpha)}, \lessdot_{f(\alpha)}}}\\
	&\text{ যদি এবং কেবল যদি }\ordeq{\tuple{A, <}}{\tuple{B_b, \lessdot_b}}\text{ কোনো $b \in B$-এর জন্য}
\end{align*}
এখানে \olref[iso]{ordisounique},
\olref[iso]{wellordinitialsegment} এবং
\olref[basic]{ordissetofsmallerord} ব্যবহার করেছি।
\end{proof}

\end{document}
